% ============================================================================
%  Capitolo 18 -- Il processore                                  (Lezione 18)
% ============================================================================
\chapter{Il processore}
\label{ch:processor}

\epigraphline{Un Anello per domarli tutti\ldots}
  {J.\,R.\,R.\ Tolkien, \emph{Il Signore degli Anelli}, 1954--1955}

\begin{chapgoals}
Alla fine di questo capitolo saprete
\begin{itemize}
\item nominare i componenti di un processore: unità di controllo, ALU, registri, collegamenti con la memoria e con l'ingresso/uscita;
\item distinguere i registri interni dai registri di macchina, e i registri generici da quelli speciali;
\item spiegare come funziona uno stack e come \riscv realizza push e pop;
\item confrontare i modelli ad accumulatore, a stack e a registri generali sullo stesso calcolo;
\item descrivere \riscv come un processore load-store a registri generali, e leggere il nome RV32IMAC;
\item far corrispondere ogni elemento del modello generale al processore del kit \prg.
\end{itemize}
\end{chapgoals}
% Generated by tools/textbook/export_listings.js from tools/decks/mastery.js and tools/concepts/registry.json: do not edit.
\begin{chapstatus}{Stato di apprendimento di questo capitolo}
\masterychip{now}{DA COMPRENDERE OGGI}\enspace accumulatore, pila, registri; il modello del processore; i registri di RISC-V; la convenzione di chiamata; la pila; il processore RV32.\par\smallskip
\masterychip{use}{DA UTILIZZARE OGGI}\enspace il programma completo, gli strumenti di PRG32 (si usa, non si studia, in questo corso).\par\smallskip
\masterychip{later}{DA COMPRENDERE PIÙ AVANTI}\enspace interruzioni e trap (Capitolo~\ref{ch:instructions}).
\end{chapstatus}


Programmate un processore fin dal Capitolo~\ref{ch:firstcontact}. Questo
capitolo presenta il modello generale di un processore, di cui ogni processore
reale è una variante, e usa \riscv come riferimento: ogni elemento astratto ha
un nome che conoscete già.

% ----------------------------------------------------------------------------
\section{Riscaldamento sulla macchina}

\asmfile[caption={Una riga di un linguaggio ad alto livello: \code{z = x + y}.},label=lst:cpu:sum]{code/l18/c03.S}

Sette righe per una somma (Listato~\ref{lst:cpu:sum}). La memoria è toccata
solo da \ins{lw} e \ins{sw}, e l'aritmetica avviene solo tra registri: è il
ciclo «carica, modifica, memorizza» del Capitolo~\ref{ch:firstcontact}, visto
ora come una scelta di progetto. Altri processori fanno lo stesso lavoro con
tre istruzioni, o con una. Perché \riscv è fatto così? Per rispondere serve un
modello del processore.

% ----------------------------------------------------------------------------
\section{Il processore}

\subsection{Il calcolatore e il suo processore}

Un calcolatore elettronico è un sistema numerico, automatico, a programma
memorizzato, costruito con circuiti elettronici (Capitolo~\ref{ch:boole1}),
con tre sottosistemi: processore, memoria centrale e ingresso/uscita,
collegati da bus. Il processore può eseguire solo un insieme fisso di
istruzioni, da alcune decine a poche centinaia, più o meno complesse. Il
\term{linguaggio macchina} è l'insieme delle stringhe di bit che codificano le
istruzioni che i circuiti di un processore possono eseguire. Le istruzioni
richiedono trasferimenti di informazione, all'interno del calcolatore e da o
verso i dispositivi di I/O, e operazioni aritmetiche e logiche. I dati e le
istruzioni di un programma sono memorizzati nella memoria centrale.

RV32I, la base di \riscv, ha circa quaranta istruzioni; con le estensioni M, A
e C del kit, circa un centinaio di codifiche. Il Motorola 68000 aveva
cinquantasei tipi di istruzione, con molte varianti di ciascuno.

\subsection{Un modello di riferimento}

La Figura~\ref{fig:cpu:model} è il modello di riferimento di un processore. I
suoi elementi sono i seguenti.

\begin{itemize}
\item La \term{rete di controllo}, o unità di controllo, che produce i comandi
  per tutto il resto.
\item Il \term{program counter} PC, che contiene l'indirizzo della prossima
  istruzione, e il \term{registro di decodifica} D, che contiene l'istruzione
  in esecuzione.
\item I \term{registri di macchina} $R_0,\dots,R_{n-1}$.
\item La \term{ALU}, con i suoi registri di ingresso $O_1$ e $O_2$ e il suo
  registro di uscita $U$.
\item Il collegamento con la memoria: il registro indirizzo MA e il registro
  buffer MB del Capitolo~\ref{ch:memory}.
\item Il collegamento con l'ingresso/uscita: un registro di ingresso RI e un
  registro di uscita RU.
\item Un \term{bus interno} che li unisce tutti.
\end{itemize}

\begin{figure}[htb]
\centering
\begin{tikzpicture}
  \node[draw=prgsteel,dashed,rounded corners,minimum width=116mm,minimum height=62mm,label={[lbl,anchor=north west]north west:CPU}] (cpu) at (5.3,-2.25) {};
  \node[blk,fill=prgsteel,text=white,minimum width=20mm,minimum height=17mm] (cn) at (1.0,-0.6) {rete di\\controllo};
  \node[blkred,minimum width=9mm,minimum height=7mm] (pc) at (0.5,-2.2) {PC};
  \node[blkred,minimum width=9mm,minimum height=7mm] (d) at (1.5,-2.2) {D};
  \foreach \i/\r in {0/{$R_0$},1/{$R_1$},2/{$\cdots$},3/{$R_{n-1}$}} \node[blkacc,minimum width=18mm,minimum height=5.5mm] (r\i) at (4.0,-\i*0.65+0.1) {\r};
  \node[blkvio,minimum width=11mm,minimum height=6mm] (o1) at (6.9,0.1) {$O_1$};
  \node[blkvio,minimum width=11mm,minimum height=6mm] (o2) at (8.3,0.1) {$O_2$};
  \node[blkvio,minimum width=25mm,minimum height=9mm] (alu) at (7.6,-0.85) {\textbf{ALU}};
  \node[blkvio,minimum width=11mm,minimum height=6mm] (u) at (7.6,-1.9) {$U$};
  \draw[prgamber,line width=3pt] (-0.2,-2.95) -- (10.8,-2.95) node[pos=0.93,above,lbl,text=prgamber]{bus interno};
  \node[blkgrn,minimum width=11mm,minimum height=7mm] (ma) at (1.0,-4.1) {MA};
  \node[blkgrn,minimum width=11mm,minimum height=7mm] (mb) at (2.6,-4.1) {MB};
  \node[blkamb,minimum width=11mm,minimum height=7mm] (ri) at (7.0,-4.1) {RI};
  \node[blkamb,minimum width=11mm,minimum height=7mm] (ru) at (8.6,-4.1) {RU};
  \foreach \n in {pc,d,u} \draw[<->,thick] (\n.south) -- (\n.south|-0,-2.95);
  \draw[<->,thick] (r3.south) -- (r3.south|-0,-2.95);
  \foreach \n in {ma,mb,ri,ru} \draw[<->,thick] (\n.north) -- (\n.north|-0,-2.95);
  \draw[->,thick] (o1) -- (o1|-alu.north); \draw[->,thick] (o2) -- (o2|-alu.north); \draw[->,thick] (alu) -- (u);
  \draw[<->,bus] (1.8,-4.45) -- (1.8,-6.0) node[below,lbl]{verso la memoria};
  \draw[<->,bus] (7.8,-4.45) -- (7.8,-6.0) node[below,lbl]{verso l'ingresso/uscita};
\end{tikzpicture}
\caption{Il modello di riferimento di un processore.}
\label{fig:cpu:model}
\end{figure}

Pur essendo basati sullo stesso schema logico, e quindi fatti di componenti
piuttosto simili, i processori in commercio possono essere molto diversi l'uno
dall'altro. La Tabella~\ref{tab:cpu:map} mette fianco a fianco il modello e
\riscv. Le due righe segnate «assente» sono semplificazioni deliberate di
\riscv, discusse più avanti e nel Capitolo~\ref{ch:memory}. L'MC68000 delle
edizioni precedenti di questo corso aveva tutti gli elementi del modello: otto
registri dati e otto registri indirizzi, un registro di stato con i codici di
condizione, e una circuiteria separata per le interruzioni.

\begin{table}[htb]
\centering\small
\caption{Il modello e \riscv.}\label{tab:cpu:map}
\begin{tabularx}{\textwidth}{@{}P{0.3\textwidth}LL@{}}
\toprule
\textbf{Elemento del modello} & \textbf{In \riscv} & \textbf{Visto dal programmatore?}\\
\midrule
PC, program counter & \reg{pc}, 32 bit & indirettamente: lo usano i salti e \ins{auipc}\\
D, registro di decodifica & il registro istruzione & no\\
$R_0\ldots R_{n-1}$, registri di macchina & \reg{x0}\ldots\reg{x31}, trentadue registri di 32 bit & sì: \reg{zero}, \reg{ra}, \reg{sp}, \reg{t0}, \reg{a0}, \reg{s0}, \ldots\\
$O_1$, $O_2$, $U$, registri della ALU & gli ingressi e l'uscita della ALU & no\\
MA, MB, collegamento con la memoria & l'indirizzo e il dato di un load o di uno store & no\\
RI, RU, collegamento con l'ingresso/uscita & assente: i dispositivi sono mappati in memoria & --\\
registro di stato con i flag & assente: nessun codice di condizione & --\\
registri speciali & i registri di controllo e stato (CSR) & solo dal software privilegiato\\
\bottomrule
\end{tabularx}
\end{table}

\subsection{Registri interni e registri di macchina}

\begin{definition}{Registri interni e di macchina}
I \term{registri interni} non sono direttamente visibili al programmatore, e
sono necessari al funzionamento del processore: PC, MA, MB, D, $O_1$, $O_2$,
$U$. I \term{registri di macchina} definiscono il modello di programmazione;
si dividono in registri \term{generici} e registri \term{speciali}.
\end{definition}

Un processore tratta tre tipi di informazione: istruzioni, indirizzi di
memoria e dati. La dimensione di MA e del PC dipende dalla dimensione degli
indirizzi; la dimensione dei registri della ALU dipende dalla dimensione dei
dati. In RV32 tutte e tre sono larghe 32 bit: una sola dimensione per
registri, indirizzi e istruzioni. Il «32» del nome è la larghezza dei
registri: RV64 ha registri e indirizzi di 64 bit, e le stesse istruzioni di 32
bit.

\subsection{In che cosa i processori reali differiscono dal modello}

In generale, i processori reali possono avere un insieme di registri generici
diverso e di solito più ricco; i loro registri di macchina possono essere
specifiche locazioni di memoria; la sezione di I/O può mancare del tutto ed
essere realizzata da chip dedicati; e unità aritmetiche molto semplici possono
essere estese da chip dedicati che migliorano la ALU, come facevano i
coprocessori in virgola mobile 68881 per la famiglia 68000 e 8087 per l'8086.

\riscv sul kit \prg ha trentadue registri generici in un banco dentro il
processore; nessuna sezione di I/O, dato che i controller delle periferiche
sono sullo stesso chip e si raggiungono come memoria; \emph{estensioni}
opzionali al posto di chip aggiuntivi, tra le quali M aggiunge moltiplicazione
e divisione; e nessuna estensione F, così che la virgola mobile andrebbe fatta
in software (Capitolo~\ref{ch:reals}).

\subsection{Parole o byte: la lunghezza delle parole di macchina}

In un \emph{processore a parola} tutte le parole di macchina hanno la stessa
lunghezza, dimensionata sull'informazione più lunga; questo è scomodo con dati
di natura intrinsecamente variabile. Ogni istruzione \riscv è lunga 32 bit (16
nella forma compressa). In un processore \emph{a carattere}, o a byte, la
lunghezza delle parole di macchina è variabile, un multiplo di un byte, il che
si adatta al trattamento delle stringhe: un'istruzione Intel x86 occupa da 1 a
15 byte, e un'istruzione MC68000 da una a cinque parole di 16 bit. Un parametro
caratteristico di un processore è il numero di bit dei suoi registri, il
parallelismo dei suoi trasferimenti interni: 8, 16, 32 o 64.

\subsection{Stato utente e stato supervisore}

Un processore è in \term{stato supervisore} quando esegue il sistema
operativo, e può allora eseguire ogni istruzione del suo insieme; è in
\term{stato utente} quando esegue un programma utente, e può allora eseguire
tutto tranne le \term{istruzioni privilegiate}.

\riscv definisce tre \term{modi di privilegio}. Il modo macchina (M), il più
privilegiato e l'unico obbligatorio, esegue il firmware. Il modo supervisore
(S) esegue il kernel di un sistema operativo come Linux. Il modo utente (U)
esegue le applicazioni. L'ESP32-C6 realizza M e U. Il passaggio da un modo
all'altro avviene attraverso le trap, argomento del
Capitolo~\ref{ch:instructions}.

% ----------------------------------------------------------------------------
\section{I registri}

\subsection{Tipi di registri}

Un registro può essere specializzato, o può fare più lavori. I registri
specializzati richiedono codici operativi specializzati per la stessa
operazione su banchi diversi; concentrare più funzioni negli stessi registri
dà maggiore complessità interna e maggiore flessibilità d'uso. I registri di
macchina, diversi in numero e spesso raggruppati in banchi, possono essere
organizzati come una memoria indirizzabile (un banco $R$, di cui $R_i$ è un
registro) o come registri autonomi, ciascuno con il proprio nome. Alcuni
processori collocavano i loro registri generici nelle prime locazioni di
memoria: la «pagina zero» del 6502.

\riscv ha un solo banco di trentadue registri, indirizzato da 5 bit dentro
l'istruzione; i nomi \reg{t0}, \reg{a0}, \reg{sp} sono solo convenzioni.

\subsection{Le funzioni dei registri}

\begin{description}[style=nextline,leftmargin=1.5em]
\item[Transito dei dati.] Dalla memoria a un registro, da registro a registro,
  da un registro alla memoria: load, move, store.
\item[Accumulo dei risultati.] Necessario per l'aritmetica e la logica:
  $R:=\mathrm{op}\,R$, oppure $R:=R\ \mathrm{op}\ X$.
\item[Indirizzamento.] Trovare l'indirizzo di un operando: registri indice,
  registri base.
\item[Puntatori allo stack.] Gestire uno stack realizzato in memoria.
\item[Indicatori, o flag.] Informazione binaria su eventi: zero, segno,
  overflow, riporto, indirizzi o codici illegali. Prodotti durante
  l'esecuzione di un'istruzione, possono condizionare le successive.
\end{description}

In \riscv le prime quattro funzioni sono tutte svolte dai registri generici;
la quinta non esiste come registro.

\subsection{I trentadue registri di RISC-V}

\begin{table}[htb]
\centering\small
\caption{I registri di \riscv e i loro ruoli convenzionali.}\label{tab:cpu:regs}
\begin{tabular}{@{}lll@{\qquad}lll@{}}
\toprule
\textbf{Registro} & \textbf{Nome} & \textbf{Funzione} & \textbf{Registro} & \textbf{Nome} & \textbf{Funzione}\\
\midrule
\reg{x0} & \reg{zero} & sempre 0 & \reg{x8}--\reg{x9} & \reg{s0}--\reg{s1} & salvati\\
\reg{x1} & \reg{ra} & indirizzo di ritorno & \reg{x10}--\reg{x11} & \reg{a0}--\reg{a1} & argomenti, risultati\\
\reg{x2} & \reg{sp} & puntatore allo stack & \reg{x12}--\reg{x17} & \reg{a2}--\reg{a7} & argomenti\\
\reg{x3} & \reg{gp} & puntatore globale & \reg{x18}--\reg{x27} & \reg{s2}--\reg{s11} & salvati\\
\reg{x4} & \reg{tp} & puntatore al thread & \reg{x28}--\reg{x31} & \reg{t3}--\reg{t6} & temporanei\\
\reg{x5}--\reg{x7} & \reg{t0}--\reg{t2} & temporanei & \reg{pc} & -- & program counter\\
\bottomrule
\end{tabular}
\end{table}

L'hardware tratta in modo identico i registri da \reg{x1} a \reg{x31}, e solo
\reg{x0} è speciale (Tabella~\ref{tab:cpu:regs}); il program counter non è nel
banco. I ruoli vengono dalla convenzione di chiamata, incontrata nel
Capitolo~\ref{ch:firstcontact}, la cui regola è questa. Una chiamata può
distruggere \reg{ra} e i registri \reg{t} e \reg{a}: sono \emph{salvati dal
chiamante}, e un chiamante che ne ha ancora bisogno li salva prima della
chiamata. Una chiamata deve restituire invariati i registri \reg{s} e \reg{sp}:
sono \emph{salvati dal chiamato}, e una funzione che ne usa uno lo salva nel suo
frame sullo stack e lo ripristina prima di \ins{ret}. L'Appendice~\ref{app:riscv}
dà la convenzione completa. L'MC68000, invece,
separava i registri dati D0--D7 dai registri indirizzi A0--A7, con A7 come
puntatore allo stack.

\subsection{I flag: il registro di stato, e la sua assenza}

Molti processori hanno un \term{registro di stato}: bit, chiamati flag,
accessibili dal programma. Quelli usuali sono Z (zero), N (negativo), V
(overflow) e C (riporto), talvolta con un bit di parità. Un'istruzione li
imposta, e un successivo salto condizionato li esamina. L'MC68000 ha il suo
registro dei codici di condizione, e x86 il suo registro dei flag.

\riscv non ha flag. Un salto confronta due registri e salta, in una sola
istruzione: \ins{beq}, \ins{bne}, \ins{blt}, \ins{bge}, \ins{bltu},
\ins{bgeu}. Le istruzioni \ins{slt} e \ins{sltu} scrivono il risultato di un
confronto in un registro generico. Riporto e overflow si calcolano quando
servono (Capitolo~\ref{ch:integers}). C'è meno stato nascosto, il che rende
più semplici le pipeline e le interruzioni.

\subsection{I registri speciali di RISC-V: i CSR}

I registri speciali di \riscv sono i \term{registri di controllo e stato}
(CSR, \emph{control and status registers}): un banco separato di un massimo di
4\,096 registri, letti e scritti solo con le istruzioni \ins{csrrw},
\ins{csrrs} e \ins{csrrc}, e solo dal codice privilegiato
(Tabella~\ref{tab:cpu:csr}). Sono i «registri delle interruzioni» del modello
generale, e l'argomento del Capitolo~\ref{ch:instructions}.

\begin{table}[htb]
\centering\small
\caption{Alcuni CSR del modo macchina.}\label{tab:cpu:csr}
\begin{tabularx}{\textwidth}{@{}llL@{}}
\toprule
\textbf{Registro} & \textbf{Nome} & \textbf{Che cosa contiene}\\
\midrule
\reg{mstatus} & machine status & l'abilitazione globale delle interruzioni e il modo di privilegio precedente\\
\reg{mtvec} & trap vector & l'indirizzo della routine che gestisce interruzioni ed eccezioni\\
\reg{mepc} & exception program counter & l'indirizzo dell'istruzione che è stata interrotta\\
\reg{mcause} & trap cause & un codice che dice perché è avvenuta la trap\\
\reg{mie}, \reg{mip} & interrupt enable, interrupt pending & un bit per sorgente di interruzione\\
\bottomrule
\end{tabularx}
\end{table}

% ----------------------------------------------------------------------------
\section{Lo stack}

\subsection{L'ultimo a entrare è il primo a uscire}

\begin{definition}{Stack}
Uno \term{stack}, o pila, è una struttura dati astratta con due operazioni:
\emph{push} memorizza un elemento come ultimo della pila, e \emph{pop} toglie
l'ultimo elemento. L'ultimo valore inserito è il primo estratto (LIFO,
\emph{last in, first out}).
\end{definition}

Uno stack lavora sempre e solo sulla sua cima: non c'è bisogno di dire dove
memorizzare o da dove prelevare. Si usa per mettere da parte il contenuto dei
registri che vanno liberati per svolgere un compito, e per ripristinarlo dopo:
accade nei sottoprogrammi e nelle interruzioni. È realizzato con un puntatore,
in memoria o in un registro dedicato, che indica in ogni istante la posizione
dell'ultimo elemento. Le operazioni su uno stack dovrebbero controllare
l'overflow su un push e lo stack vuoto su un pop.

\subsection{Push e pop in RISC-V}

In \riscv lo stack vive in memoria, cresce verso gli indirizzi bassi, e il
registro \reg{sp} punta alla sua cima (Figura~\ref{fig:cpu:stack}). Non ci sono
istruzioni chiamate push e pop: ciascuna è fatta di due istruzioni ordinarie.

\begin{figure}[htb]
\centering
\begin{minipage}{0.42\textwidth}\centering
\begin{tikzpicture}[x=0.7cm,y=0.62cm,font=\ttfamily\small]
  \foreach \i/\v/\f in {0/(libero)/prgpanel, 1/(libero)/prgpanel, 2/-28/prgamber!60, 3/-17/prgaccent!18, 4/739/prgaccent!18, 5/43/prgaccent!18}{
    \draw[prgsteel,fill=\f] (0,-\i) rectangle (2.6,-\i-1);
    \node at (1.3,-\i-0.5) {\v};
  }
  \draw[->,very thick,prgred] (-1.1,-2.5) node[left,font=\ttfamily\bfseries]{sp} -- (-0.05,-2.5);
  \node[lbl,anchor=west,text=prgcomment] at (2.7,-0.5) {indirizzi bassi};
  \node[lbl,anchor=west,text=prgcomment] at (2.7,-5.5) {fondo};
\end{tikzpicture}
\end{minipage}
\begin{minipage}{0.5\textwidth}
\begin{asm}[numbers=none,xleftmargin=4pt,basicstyle=\ttfamily\footnotesize]
# push t0
addi sp, sp, -4      # make room
sw   t0, 0(sp)       # store on the new top

# pop t0
lw   t0, 0(sp)       # read the top
addi sp, sp, 4       # release it
\end{asm}
\end{minipage}
\caption{Uno stack in memoria, e le due operazioni in \riscv.}
\label{fig:cpu:stack}
\end{figure}

In pratica un programma \riscv non inserisce ed estrae valori singoli: uno
\emph{stack frame} è allocato una volta per funzione, sedici byte alla volta,
ed è riempito con store ordinari, come ha fatto ogni funzione di \vermone che
fa una chiamata. Altri processori hanno forme dedicate: sul 68000,
\code{MOVE (SP)+, D0} estrae con un post-incremento.

% ----------------------------------------------------------------------------
\section{Modelli di processore}

I processori si classificano secondo il luogo in cui vivono gli operandi di
un'operazione.

\subsection{Quattro architetture}

\begin{description}[style=nextline,leftmargin=1.5em]
\item[A registri generali.] Il processore ha un insieme di registri che
  possono essere usati in modo intercambiabile.
\item[Ad accumulatore.] Tutte le istruzioni aritmetiche, logiche e di
  confronto hanno un operando implicito: un registro speciale, l'accumulatore.
\item[A stack.] Non ci sono né accumulatori né registri generali, solo il
  puntatore allo stack; le istruzioni lavorano sui dati in cima allo stack.
\item[Ortogonale.] Un solo banco di registri generali, alcuni dei quali
  servono, quando occorre, da puntatore allo stack o da altri registri
  speciali.
\end{description}

\subsection{\texttt{z = x + y} in quattro processori}

\begin{table}[htb]
\centering\small
\caption{Lo stesso calcolo in quattro modelli.}\label{tab:cpu:models}
\begin{tabularx}{\textwidth}{@{}P{0.21\textwidth}P{0.2\textwidth}Lc@{}}
\toprule
\textbf{Modello} & \textbf{Esempio} & \textbf{Programma} & \textbf{Accessi}\\
\midrule
accumulatore & Motorola 6809, MOS 6502 & \code{LDA x; ADD y; STA z} & 3\\
stack & macchina virtuale Java, WebAssembly & \code{PUSH x; PUSH y; ADD; POP z} & 3 + stack\\
registri generali, operandi in memoria & MC68000, Intel x86 & \code{MOVE x,D0; ADD y,D0; MOVE D0,z} & 3\\
registri generali, load-store & \riscv & \code{lw t1,x; lw t2,y; add t3,t1,t2; sw t3,z} & 3\\
\bottomrule
\end{tabularx}
\end{table}

Gli accessi alla memoria sono dappertutto gli stessi tre
(Tabella~\ref{tab:cpu:models}): ciò che cambia è quante istruzioni ci sono, e
quanto è complessa ciascuna. Nel modello registro-memoria un'istruzione come
\code{ADD y, D0} legge la memoria \emph{e} calcola; nel modello load-store
ogni istruzione fa una cosa sola. Le macchine a stack con lo stack in memoria
sono generalmente più lente delle macchine a registri generali. (Nella
tabella, \code{lw t1, x} è un'abbreviazione che l'assemblatore accetta per la
coppia \ins{la} e \ins{lw} del riscaldamento.)

\subsection{RISC-V: un processore load-store a registri generali}

\begin{prgconcept}[Caricare, calcolare, memorizzare]
Solo i load e gli store toccano la memoria; ogni altra istruzione lavora da
registro a registro. È il ritmo di ogni programma \riscv, dal primo contatore
all'ultimo gioco.
\end{prgconcept}

Le ragioni sono quattro. Le istruzioni che lavorano sui registri sono più
veloci di quelle che lavorano su locazioni di memoria, così che i dati più
usati si tengono nei registri. Un'istruzione ha tre operandi, due sorgenti e
una destinazione, così che nessun valore viene distrutto a meno che il
programmatore non lo voglia. Ogni istruzione fa una cosa semplice e ha una
dimensione fissa: è facile da decodificare, e facile da eseguire in una
pipeline. È questo il significato di RISC, \emph{reduced instruction set
computer}: le forme che vanno dalla memoria a un registro e da un registro
alla memoria, che esistono nei processori CISC, sono esattamente ciò a cui il
RISC rinuncia.

% ----------------------------------------------------------------------------
\section{La ALU e il quadro d'insieme}

\subsection{L'unità aritmetico-logica}

La ALU del Capitolo~\ref{ch:combinational}, vista dal processore, sovrintende
alle operazioni aritmetiche e logiche tra operandi. Nel modello ha due
registri di ingresso, $O_1$ e $O_2$, e un registro di uscita $U$. È
fisicamente inclusa nel processore nelle architetture elementari ed è un
apparato indipendente in quelle più complesse; una soluzione diffusa è
un'unità aritmetica ridotta dentro il processore con apparati indipendenti per
le operazioni più sofisticate. In \riscv quell'apparato sono le estensioni: M
per moltiplicazione e divisione, F e D per la virgola mobile. Le operazioni
disponibili dipendono dai tipi di dati e dalla loro rappresentazione: valori
logici, stringhe di caratteri, puntatori, e numeri nelle codifiche dei
Capitoli \ref{ch:integers} e~\ref{ch:reals}.

\subsection{L'architettura di un processore RV32}

La Figura~\ref{fig:cpu:datapath} riduce ai suoi blocchi il classico datapath
\riscv a ciclo singolo di \citet{patterson2020}. Il PC dà l'indirizzo
dell'istruzione; l'istruzione viene prelevata dalla memoria; il decodificatore
ne separa i campi e l'unità di controllo produce i comandi; il banco dei
registri fornisce gli operandi; la ALU calcola un risultato o forma un
indirizzo; la memoria dati viene letta o scritta; e il risultato torna al
banco dei registri. La memoria istruzioni e la memoria dati sono lo stesso
spazio degli indirizzi, disegnato due volte per chiarezza.

\begin{figure}[htb]
\centering
\begin{tikzpicture}[node distance=4mm]
  \node[blkred,text width=12mm,minimum height=16mm] (pc) {\textbf{PC}};
  \node[blkvio,text width=19mm,minimum height=16mm,right=of pc] (im) {\textbf{memoria istruzioni}\\\scriptsize prelievo};
  \node[blk,fill=prgsteel,text=white,text width=20mm,minimum height=16mm,right=of im] (dec) {\textbf{decodifica e controllo}\\\scriptsize campi, comandi};
  \node[blkacc,text width=19mm,minimum height=16mm,right=of dec] (rf) {\textbf{banco dei registri}\\\scriptsize \texttt{x0}--\texttt{x31}};
  \node[blkamb,text width=16mm,minimum height=16mm,right=of rf] (alu) {\textbf{ALU}\\\scriptsize calcola, o forma un indirizzo};
  \node[blkvio,text width=16mm,minimum height=16mm,right=of alu] (dm) {\textbf{memoria dati}\\\scriptsize load, store};
  \foreach \a/\b in {pc/im,im/dec,dec/rf,rf/alu,alu/dm} \draw[->,bus] (\a) -- (\b);
  \draw[->,bus,prgred,rounded corners=2mm] (dm.south) -- ++(0,-8mm) -| (rf.south) node[pos=0.25,below,lbl,text=prgred]{riscrittura del risultato};
  \draw[->,bus,prgsteel,rounded corners=2mm] (alu.north) -- ++(0,8mm) -| (pc.north) node[pos=0.25,above,lbl]{PC successivo: $PC+4$ o la destinazione di un salto};
\end{tikzpicture}
\caption{L'architettura di un processore RV32, ridotta ai suoi blocchi.}
\label{fig:cpu:datapath}
\end{figure}

\subsection{Il processore del kit PRG32}

\begin{table}[htb]
\centering\small
\caption{L'ESP32-C6 (dal datasheet di Espressif).}\label{tab:cpu:c6}
\begin{tabularx}{\textwidth}{@{}lL@{}}
\toprule
\textbf{Caratteristica} & \textbf{ESP32-C6}\\
\midrule
Insieme di istruzioni & RV32IMAC: insieme intero di base, moltiplicazione e divisione, operazioni atomiche, istruzioni compresse\\
Core principale & un core \riscv a 32 bit, fino a 160 MHz, pipeline a quattro stadi\\
Secondo core & un core \riscv a basso consumo, fino a 20 MHz, per i compiti che girano mentre il core principale dorme\\
Modi di privilegio & macchina e utente\\
Virgola mobile & nessuna: né estensione F né D\\
Memoria & 512 KB di SRAM, 320 KB di ROM, flash esterna\\
Sullo stesso chip & radio Wi-Fi 6, Bluetooth, Thread e Zigbee; periferiche SPI, I2S, GPIO e timer\\
\bottomrule
\end{tabularx}
\end{table}

Il chip della Tabella~\ref{tab:cpu:c6} è un calcolatore completo su pochi
millimetri quadrati che costa circa due euro: processore, memoria e
ingresso/uscita ci sono tutti, come nel modello con cui il capitolo è
cominciato.

% ----------------------------------------------------------------------------
% Generated by tools/textbook/export_listings.js from decks/plan.js and code/l18.S: do not edit.
\section{Il programma di questo capitolo}\label{sec:lab:l18}

I listati di questo capitolo sono parti di un solo programma,
\fname{code/l18.S} nel repository del corso, stampato per intero nel
Listato~\ref{lst:lab:l18}. I test automatici del materiale del corso lo
compilano e lo eseguono sull'esecutore senza interfaccia di \prgqt{}, con
ingressi programmati, e controllano il frame; la
Figura~\ref{fig:lab:l18} è il frame scritto dal test.

\begin{description}[style=nextline,leftmargin=1.5em]
\item[Sullo schermo.] Tre barre, otto pixel per unità: x = 5 (bianca), y = 7 (bianca) e z = 12 (verde), calcolata dalle sette istruzioni della slide.
\item[Che cosa premere.] Niente: il programma va da solo.
\end{description}

\begin{figure}[htb]
\centering
\includegraphics[width=0.5\textwidth]{lab_l18.png}
\caption{\fname{l18.S} su \prgqt{}.}\label{fig:lab:l18}
\end{figure}

\asmfile[caption={Il programma completo del capitolo, \fname{code/l18.S}.},label=lst:lab:l18]{code/l18/program.S}

Per eseguirlo su \prgqt{} i comandi sono quelli del
Listato~\ref{lst:lab:l18:run}; il Paragrafo~\ref{sec:prg:testing} dà gli
stessi passi per QEMU e per la scheda.

\begin{shell}[caption={Compilare \fname{l18.S} ed eseguirlo su \prgqt{}.},label=lst:lab:l18:run]
# 1. start PRG32-QT on this computer (it listens on port 8080)

# 2. in the PRG32 folder, with ESP-IDF active: build the cartridge
python3 -m prg32 cartridge build ~/arc12/code/l18.S --portable \
    --entry-prefix l18 --name l18 --out l18.prg32

# 3. send it to PRG32-QT and start it
python3 -m prg32 esp32c6 upload-and-run l18.prg32 \
    --url http://localhost:8080 --slot cart0

# 4. look, press, compare with the previous slide
\end{shell}


% ----------------------------------------------------------------------------
% Generated by tools/textbook/export_listings.js from tools/concepts/assessments.json: do not edit.
\begin{chapstatus}{Il filo del capitolo}
Il ra salvato nella Lezione 02 è ora spiegato: lo stack e la convenzione di chiamata sono regole che permettono a funzioni scritte da persone diverse di chiamarsi. Prossimo passo: come un'istruzione è codificata ed eseguita, e come il processore viene interrotto.
\par\medskip\textbf{Biglietto d'uscita} (da due a quattro minuti, da soli, su carta):
\begin{enumerate}[nosep,leftmargin=2em]
\item Seguite add t2, t0, t1 nello schema a blocchi di un processore RV32: nominate, in ordine, i blocchi che usa e che cosa fa ciascuno. Quale blocco non usa?
\item Una funzione comincia con addi sp, sp, -16 e sw ra, 12(sp). Dove si trova ora ra, e che cosa deve fare la funzione prima di ret?
\end{enumerate}
\end{chapstatus}


\section{Verificate di aver capito}

\begin{quiz}{Quale registro del modello generale non ha corrispettivo in \riscv?}
\item il program counter
\item il registro di stato con i flag
\item il puntatore allo stack
\item il registro istruzione
\end{quiz}

\begin{quiz}{In \riscv, quali istruzioni possono leggere o scrivere la memoria centrale?}
\item tutte
\item \ins{add} e \ins{sub}
\item solo load e store
\item solo i salti
\end{quiz}

\begin{quiz}{Dopo \code{addi sp, sp, -4} e \code{sw t0, 0(sp)}, che cosa è successo?}
\item \reg{t0} è stato estratto dallo stack
\item \reg{t0} è stato inserito sullo stack
\item lo stack è stato svuotato
\item \reg{sp} è stato salvato in \reg{t0}
\end{quiz}

\begin{quiz}{In una macchina ad accumulatore, l'istruzione \code{ADD y} significa:}
\item $y := y + \text{accumulatore}$
\item $\text{accumulatore} := \text{accumulatore} + [y]$
\item push di $y$, poi somma
\item $y := \text{accumulatore}$
\end{quiz}

\section{Esercizi}

\begin{exercise}[prima su carta, poi su \prg]
\begin{enumerate}[nosep]
\item Scrivete $z=(a+b)-(c+d)$ per una macchina ad accumulatore, per una
  macchina a stack e per \riscv; contate le istruzioni e gli accessi alla
  memoria.
\item Scrivete il push e il pop di due registri, \reg{t0} e \reg{t1}, tenendo
  \reg{sp} multiplo di 16 come richiede la convenzione di chiamata.
\item Per ogni registro usato nel vostro ultimo programma \prg, dite quale
  funzione svolge: transito, accumulo, indirizzamento, stack.
\item In \prgqt, eseguite passo-passo una chiamata e il suo ritorno: guardate
  \reg{pc}, \reg{ra} e \reg{sp}, e la parola salvata sullo stack.
\item Trovate nel datasheet del processore del vostro telefono o portatile:
  insieme di istruzioni, numero di core, clock, livelli di privilegio.
\end{enumerate}
\expected{Accumulatore: \code{LDA a; ADD b; STA tmp; LDA c; ADD d; STA tmp2;
LDA tmp; SUB tmp2; STA z}. Stack: \code{PUSH a; PUSH b; ADD; PUSH c; PUSH d;
ADD; SUB; POP z}. \riscv: quattro load, due \ins{add}, una \ins{sub}, uno
store. In \prgqt: \reg{sp} scende di 16 al prologo e ritorna all'epilogo;
\reg{ra} contiene l'indirizzo dopo la chiamata.}
\end{exercise}

\section{Glossario}

\begin{glossarytable}
control unit & unità di controllo\\
program counter & contatore di programma\\
machine register & registro di macchina\\
general / special register & registro generico / speciale\\
instruction set & insieme (set) di istruzioni\\
accumulator & accumulatore\\
stack pointer & puntatore allo stack\\
status register, flags & registro di stato, indicatori\\
supervisor / user state & stato supervisore / utente\\
privileged instruction & istruzione privilegiata\\
\end{glossarytable}

\section{Conclusioni e letture}

Il processore comanda e controlla, appoggiandosi alla memoria e ai sistemi di
ingresso/uscita. I registri sono locazioni di memorizzazione che possono
trovarsi fisicamente dentro o fuori il processore, e i tipi di registri e il
modo in cui sono usati determinano il modello del processore: accumulatore,
stack, registri generali. La ALU svolge le operazioni tra operandi. \riscv è
un processore load-store a registri generali, senza flag e senza istruzioni di
I/O: pochi elementi, ciascuno semplice.

\paragraph{Letture.} \citet{patterson2020}, capitoli 2 e~4; il manuale
dell'insieme di istruzioni \riscv, volume~I \citep{riscvisa};
\citet{stallings2021}, i capitoli sulla struttura e sul funzionamento del
processore. L'Appendice~\ref{app:riscv} è il riferimento completo.
